% !PRINTED 87
\weberpage{87}

\begin{center}
{\large \S.\ 25.}

\vspace{0.5em}
{\large\gesperrt{Elimination aus linearen Gleichungen.}}
\end{center}

Es kommt bisweilen vor, dass es sich bei einem gegebenen System linearer Gleichungen nicht sowohl um die wirkliche Ermittelung der Unbekannten handelt, als um die Beurtheilung der Möglichkeit ihrer Lösung, also um die Aufstellung der Bedingungsgleichungen, die zwischen den Coëfficienten bestehen müssen, wenn Lösungen oder Lösungen von bestimmter Art überhaupt vorhanden sein sollen. Die Aufstellung dieser Bedingungsgleichungen heisst \gesperrt{Elimination}. Implicite ist die Lösung dieser Aufgabe schon im Vorhergehenden enthalten; wir wollen aber noch ausdrücklich auf einige hierher gehörige Fragen zurückkommen.

Wir betrachten, wie im vorigen Paragraphen, ein System von $m$ linearen Gleichungen mit $n$ homogen vorkommenden Unbekannten, und fragen: wann hat dies System eine Lösung, bei der nicht alle Unbekannten verschwinden? Wir haben schon gesehen, dass dies immer der Fall ist, wenn $n>m$ ist. Ist aber $n\le m$, so ist die \gesperrt{nothwendige und hinreichende} Bedingung für eine solche Lösung die, dass alle $n$-reihigen Determinanten der Matrix verschwinden. Denn wenn eine von diesen nicht verschwindet, so sind nach \S.\,24, II. die Werthe der Unbekannten nothwendig Null, während, wenn sie alle verschwinden, eine Zahl $\nu<n$ gefunden werden kann, so dass alle $(\nu+1)$reihigen Determinanten der Matrix Null sind, während von den $\nu$-reihigen wenigstens eine nicht verschwindet, so dass also nach \S.\,24, III eine Lösung von der verlangten Art vorhanden ist.

Nun lassen sich, wenn $n\le m$ ist, aus der Matrix \S.\,24, (3)
\[
 {m(m-1)\cdots(m-n+1)\over 1\cdot2\cdots n}
\]
$n$-reihige Determinanten bilden, und so gross wäre also die Anzahl der Bedingungen. Ist $n=m$, so ist diese Zahl $=1$ und wir erhalten den Fall \S.\,24, II. und wie zu erwarten war, eine Bedingung. Im Allgemeinen ist aber diese Anzahl der Bedingungen, obwohl sie alle erfüllt sein müssen, grösser als nöthig ist, weil einige von ihnen nothwendige Folgen der übrigen sind.
